% Corpo docente e distribuição por macroárea do DEE

\section{Corpo docente}\label{sec:corpo}

O DEE reúne \textbf{34 professores}, com atuação distribuída entre as principais macroáreas do Departamento. Essa composição permite integrar ensino, pesquisa aplicada, extensão e desenvolvimento tecnológico.

\begin{center}
\begin{tcolorbox}[
  colback=white,
  colframe=azulpetroleo!35,
  boxrule=0.6pt,
  arc=1.5mm,
  left=3mm,right=3mm,top=3mm,bottom=3mm,
  width=\textwidth
]
\begin{minipage}[t]{0.28\textwidth}
\centering
{\sffamily\fontsize{28}{30}\selectfont\bfseries\color{azulpetroleo} 34}\\[1mm]
{\sffamily\large\bfseries docentes}\\[1mm]
{\small\color{cinzamedio} composição total do Departamento}
\end{minipage}
\hfill
\begin{minipage}[t]{0.66\textwidth}
\small
\textbf{Atuações docentes por macroárea}\\[-0.2mm]
{\footnotesize\color{cinzamedio}Respostas com sobreposição entre áreas}\\[1.5mm]

\noindent\textbf{Automação} \hfill \textbf{10}
\begin{tikzpicture}[baseline=(current bounding box.center)]
\fill[azulpetroleo!15] (0,0) rectangle (10,0.30);
\fill[azulpetroleo] (0,0) rectangle (7.1,0.30);
\end{tikzpicture}\\[1.5mm]

\noindent\textbf{Eletrônica e Telecomunicações} \hfill \textbf{14}
\begin{tikzpicture}[baseline=(current bounding box.center)]
\fill[azulpetroleo!15] (0,0) rectangle (10,0.30);
\fill[azulpetroleo] (0,0) rectangle (10,0.30);
\end{tikzpicture}\\[1.5mm]

\noindent\textbf{Eletrotécnica} \hfill \textbf{14}
\begin{tikzpicture}[baseline=(current bounding box.center)]
\fill[azulpetroleo!15] (0,0) rectangle (10,0.30);
\fill[azulpetroleo] (0,0) rectangle (10,0.30);
\end{tikzpicture}\\[1.5mm]

\noindent\textbf{Sistemas de Energia} \hfill \textbf{6}
\begin{tikzpicture}[baseline=(current bounding box.center)]
\fill[azulpetroleo!15] (0,0) rectangle (10,0.30);
\fill[azulpetroleo] (0,0) rectangle (4.3,0.30);
\end{tikzpicture}
\end{minipage}
\end{tcolorbox}
\end{center}

{\small\color{cinzamedio}Alguns docentes atuam em mais de uma área; por isso, os quantitativos por macroárea não são exclusivos e sua soma é maior que o total de 34 professores.}

A composição multidisciplinar do corpo docente permite integrar competências tradicionais da eletroeletrônica a tecnologias digitais emergentes em atividades de ensino, pesquisa, extensão e P\&D.
